% Appendix G — Data and methods, full description

\paragraph{Repository universe.} \PLACEHOLDERPRACTIVE{} critical OS software repositories selected by OpenSSF \texttt{criticality\_score}~\citep{ossf-criticality} (default\_score, snapshot \PLACEHOLDERSNAPSHOT{}), enriched via the GitHub REST API and filtered to drop forks, archived, disabled, and repos not pushed within the analysis window. Non-goal: GitHub-wide representativeness. Full selection procedure in \Cref{app:repo-selection}.

\paragraph{PR universe.} All PRs \emph{created} between 2025-04-01 and 2026-03-31 in those repositories (\PLACEHOLDERNPRS{} PRs), fetched via GitHub GraphQL with \texttt{orderBy:{CREATED\_AT,DESC}}. We record metadata, commits, reviews, review-thread comments, issue comments, and merge/close timeline events.

\paragraph{Event-level AI classification.} Adapting \citet{ghaleb2026fingerprinting}'s heuristics to event granularity, we classify each actor--event pair using two signals: the actor's GitHub login (matched against our 20-entry AI bot allowlist: \texttt{copilot-swe-agent[bot]}, \texttt{devin-ai-integration[bot]}, \texttt{cursor-agent}, \texttt{claude[bot]}, \texttt{coderabbitai}, \texttt{cubic-dev-ai}, \texttt{gemini-code-assist}, \texttt{copilot-pull-request-reviewer}, etc.\ -- see \Cref{app:allowlist}), and the event's payload text (commit message or PR body), scanned for a \texttt{Co-Authored-By} trailer naming an AI tool. We use a three-way taxonomy:
\begin{itemize}
\item \textbf{AI bot} --- actor login is on the AI bot allowlist.
\item \textbf{AI powered} --- actor login is human, but the payload carries an AI \texttt{Co-Authored-By} trailer.
\item \textbf{Human} --- neither of the above. This bucket may include silent AI support whose use was not declared via a trailer or bot login.
\end{itemize}
Pooled, AI bot $+$ AI powered events are referred to as \textbf{AI authored}. Handle/name mentions in free text (e.g.\ ``\texttt{@claude please fix}'') are \emph{not} treated as AI evidence: a human typing ``\texttt{@claude}'' in a comment is not themselves an AI contributor. Detected non-AI/maintenance bot events (Dependabot, Renovate, CI bots, etc.) are tracked separately and dropped from the AI-vs-human comparisons.

\paragraph{Chain length (RQ1).} Within a PR, order attributed events by timestamp. An \emph{AI$\to$AI chain} is a maximal contiguous subsequence of events with \texttt{actor\_type=AI bot}. AI powered events do not extend a chain --- chains track successive bot turns, not authorship attribution --- so this metric is a strict subset of AI authored. The per-PR metric is the longest such chain. We report the monthly share of PRs with chain $\geq$1, $\geq$2, $\geq$5.

\paragraph{AI verdict extraction.} AI bots issue native \texttt{APPROVED}/\texttt{CHANGES\_REQUESTED} review states on only \PLACEHOLDERAIOPRSNATIVE{} PRs in our window. The remaining ${\sim}70{,}000$ AI bot reviews carry the \texttt{COMMENTED} state but embed a structured verdict line that the bot itself writes in a stable template (e.g.,\ CodeRabbit's \texttt{Actionable comments posted: N}, Cubic's \texttt{N issues found}, Cursor Bugbot's \texttt{found N potential issues}, Sourcery's intro line). We extract verdicts via per-bot regex on these structured headers --- never on free-text sentiment. The full pattern catalogue and corpus-wide yield (\PLACEHOLDERAIOPPRS{} unique PRs with $\geq 1$ AI opinion after parsing, vs.\ \PLACEHOLDERAIOPRSNATIVE{} with native explicit votes only) is in \Cref{app:ai-verdict-regex}. Gemini Code Assist reviews are free-prose without a stable structured header and are left unclassified.

\paragraph{AI vs.\ human judge counters (RQ2).} Unlike the RQ3 DiD, which pools native and regex-parsed AI verdicts on the Dual-review cohort (Any AI \& human), the RQ2 ``explicit AI review'' and ``AI-only review'' monthly shares count native \texttt{APPROVED}/\texttt{CHANGES\_REQUESTED} review states only and impose no commit-ordering constraint between the AI and human reviews, so they are a strict native-only lower bound on AI's share as judge.

\paragraph{Baseline approval rates.} As a reference point: across all PRs, human-authored PRs are approved (merged with $\geq 1$ approving review) at 50.0\% ($n$=2,606,804), versus 37.3\% for AI-authored PRs ($n$=81,341).

\paragraph{Within-PR AI-AI bias DiD (RQ 3).} For each PR we assign \texttt{author\_type}~$\in\{$AI authored, human$\}$. ``AI authored'' here is the rolled-up category --- it covers both AI bot author logins and human-author commits/PR bodies that carry an AI co-author trailer (i.e.\ AI powered). Commits authored after the PR was opened are excluded so review-cycle AI activity does not retroactively mark the original PR as AI authored. The DiD universe is the \emph{Dual-review cohort (Any AI \& human)}: PRs that received both an AI bot opinion (native or parsed) AND a human explicit review, with no commits pushed strictly between the AI's first opinion timestamp and the human's first explicit-review timestamp (so both reviewers evaluated the same commit graph). We compute the four reviewer-type $\times$ author-type cells with Wilson 95\% CIs and fit a logistic regression with HC0 standard errors on the long-format dataset \texttt{Pr(approve) $\sim$ author\_AI + reviewer\_AI + author\_AI:reviewer\_AI}; the interaction coefficient is the DiD estimate. Anchoring is tested as a robustness check by stratifying on whether an AI bot has reviewed (\Cref{app:anchoring}).

\paragraph{Within-PR within-family DiD (RQ 3, Claude check).} The pooled DiD treats all AI bots as one ``AI'' reviewer cell. To test whether \emph{Claude-based} reviewers prefer Claude-authored code more than humans do, we re-run the same logit-DiD framework on the \emph{Dual-review cohort (Claude \& human)}: PRs with a Claude opinion AND a human explicit review (with the same no-commits-between filter as the pooled cohort); authors restricted to Claude vs human; Claude reviewer approves vs human reviewer approves. ``Claude opinion'' for this analysis combines native \texttt{APPROVED}/\texttt{CHANGES\_REQUESTED} states with regex-parsed verdict-line headers from Claude's review bodies and PR-thread \texttt{issue\_comment}s (\texttt{**Recommendation**: Approve}, \texttt{**Recommendation**: Request changes}, etc. --- \Cref{app:ai-verdict-regex}). Hand-audit precision on the parsed reject pattern is 122/122 = 100\% across the corpus, with zero false positives in a 37,506-comment scan of non-Claude logins.

\paragraph{Scope: authors and reviewers, not mergers.} We measure AI participation through authoring and reviewing, not merging. No merges in our universe are attributed to AI bot accounts on our allowlist, so ``AI merger'' is not an informative dimension in this window. We discuss the limitations of that framing---including the fact that AI can still act \emph{through} a human merger---in §\ref{sec:focus}.
